% Generated draft: requires mathematical and source review.
\documentclass[11pt]{article}
\usepackage[margin=25mm]{geometry}
\usepackage{fontspec}
\setmainfont{Latin Modern Roman}
\usepackage{amsmath,amssymb,mathtools,unicode-math}
\setmathfont{Latin Modern Math}
\usepackage{enumitem,needspace,xurl,hyperref}
\setlength{\parindent}{0pt}
\setlength{\parskip}{5pt}
\setlength{\emergencystretch}{3em}
\newcommand{\N}{\mathbb{N}}
\newcommand{\Z}{\mathbb{Z}}
\newcommand{\Q}{\mathbb{Q}}
\newcommand{\R}{\mathbb{R}}
\newcommand{\C}{\mathbb{C}}
\newcommand{\F}{\mathbb{F}}
\newcommand{\A}{\mathbb{A}}
\newcommand{\PP}{\mathbb{P}}
\newcommand{\E}{\mathbb{E}}
\newcommand{\eps}{\varepsilon}
\newcommand{\dd}{\,\mathrm{d}}
\begin{document}
\section*{UnsolvedMath ID: 30000366}
\textbf{Source problem code:} catalogue entry
\subsection{Definitions and mathematical statement}
Let \(X\) be a smooth projective complex K3 surface, and let \(D^b(X)=D^b(\operatorname{Coh}X)\). Write \(\operatorname{Aut}^{\mathrm{ex}}(D^b(X))\) for the group of \(\C\)-linear exact autoequivalences up to natural isomorphism. The Mukai lattice is
\[
\widetilde H(X,\Z)=H^0(X,\Z)\oplus H^2(X,\Z)\oplus H^4(X,\Z),
\quad
\langle(r,c,s),(r',c',s')\rangle=c\cdot c'-rs'-r's.
\]
Its weight-two Hodge structure uses \(\widetilde H^{2,0}=H^{2,0}(X)\), with \(H^0\) and \(H^4\) of type \((1,1)\). An integral Hodge isometry preserves this pairing, integral lattice, and Hodge decomposition. The natural cohomological action gives
\[
\rho:\operatorname{Aut}^{\mathrm{ex}}(D^b(X))\longrightarrow
O_{\mathrm{Hdg}}(\widetilde H(X,\Z)).
\]
The general task is to determine \(\operatorname{im}\rho\) for every algebraic K3 surface \(X\).

The particular question asks whether the stated orientation-reversing Hodge isometry acting by \(-\operatorname{id}\) on \(H^2(X)\) belongs to this image. Orientation here is the standard orientation of positive four-dimensional subspaces of the real Mukai lattice, which has signature \((4,20)\). A standard specific isometry with the stated behavior is
\[
\iota(r,c,s)=(r,-c,s).
\]
Its action on \(H^0\oplus H^4\) is the identity. The source does not explicitly give that action, so the display is the conventional intended representative requiring confirmation, rather than a claim that restriction to \(H^2\) uniquely specifies an isometry. In particular, \(-\operatorname{id}\) on the whole Mukai lattice is a different isometry and must not replace the requested orientation-reversing one.

The supplied assessment checked 2026-08-21 reports the cohomological action described, and full autoequivalence-group results for generic and Picard-rank-one K3 surfaces, while broader full-group questions remain conjectural. Determining the image above is distinct from determining its kernel or the entire autoequivalence group. The partially solved label is therefore broader than, and does not on its own establish the status of, the exact image and membership questions. The precise image theorem and orientation conclusion are not reproduced in the assessment and must be checked in the supplied reference.
\subsection{Short English statement}
Determine exactly which integral Hodge isometries of a projective K3 Mukai lattice are induced by derived autoequivalences, especially the stated orientation-reversing involution. The image question is distinct from the broader unresolved classification of the whole autoequivalence group.
\subsection{Sources}
\begin{itemize}
\item Local UnsolvedMath registry, permanent ID 30000366, catalogue entry, \emph{Hodge Isometries from K3 Derived Autoequivalences}.
\item \emph{Mini-Workshop: Heterotic Strings, Derived Categories, and Stacks} (2006), Oberwolfach Reports, pp. 3019--3060, DOI \url{https://doi.org/10.4171/owr/2005/53}.
\item Supplied reference ``derived equivalences and orientation'': \url{https://arxiv.org/abs/0710.1645}.
\item Supplied reference ``derived automorphism groups of Picard-rank-one K3 surfaces'': \url{https://arxiv.org/abs/1507.06435}.
\item Supplied literature triage checked 2026-08-21; it distinguishes cohomological-action results from full-group classification, without reproducing exact theorem statements.
\end{itemize}
\end{document}
